\section{Related Work} \label{sec:related}

This section surveys existing research across four relevant areas: cyber-physical system (CPS) attack analysis, fault-injection frameworks, intrusion detection for autonomous systems, and reproducible evaluation benchmarks.
For each area we summarise the most relevant contributions, highlight their limitations relative to our goals, and conclude with a structured comparison.

\subsection{CPS Attack Analysis}

Shoukry et al.~\cite{shoukry2017safecontrol} formalised secure state estimation under sensor attacks using satisfiability modulo theory, demonstrating that GPS spoofing can defeat Kalman-filter-based localisation in ground vehicles.
Their work provides a theoretical foundation but evaluates only a single attack type (sensor corruption) against a single estimator, with no orchestration or defence comparison.

Mo and Sinopoli~\cite{mo2009replay} analysed replay attacks on linear dynamical systems and proposed a $\chi^2$-detector countermeasure.
While foundational, the evaluation is purely analytical and lacks implementation on actual (or simulated) autonomous platforms.

Quarta et al.~\cite{quarta2017industrialrobot} performed a practical security analysis of an industrial robot controller, demonstrating command injection and firmware manipulation.
Their approach is platform-specific (ABB IRB 140) and does not generalise to multi-platform or multi-attack evaluation.

Mekdad et al.~\cite{mekdad2023uavsecurity} provide a broad survey of UAV security and privacy issues, systematically classifying threats at four levels: hardware, software, communication, and sensor.
The survey catalogues vulnerabilities and countermeasures but does not provide an evaluation framework or runnable code.

Cui and Stolfo~\cite{cui2010embedded} performed a wide-area scan of embedded network devices and quantified the prevalence of default credentials and unpatched firmware in deployed systems.
Their results confirm the scale of embedded-device insecurity but focus on network-accessible devices rather than protocol-level attacks on autonomous platforms.

Sun et al.~\cite{sun2020lidar} demonstrated a black-box adversarial sensor attack against LiDAR-based perception in autonomous driving, achieving approximately 80\% success rate across target models.
They proposed countermeasures (CARLO and SVF) but evaluated only a single sensor modality without multi-attack orchestration.

\textbf{Limitation:} These works evaluate individual attack vectors in isolation without providing reusable infrastructure for systematic multi-attack comparison.

\subsection{Fault-Injection Frameworks}

Jha et al.~\cite{jha2019drivefi} proposed DriveFI, a machine-learning-based fault-injection engine for autonomous vehicles that uses Bayesian reasoning to identify safety-critical faults.
DriveFI found 561 safety-critical faults in under four hours on NVIDIA and Baidu AV stacks, but it targets sensor and actuator faults rather than \emph{adversarial protocol-level attacks}, and includes no defence evaluation mechanism.

Natella et al.~\cite{natella2016faultinjection} provided a broad survey of software fault injection techniques, establishing a taxonomy of injection points, fault models, and assessment criteria.
Their contribution is taxonomic. No executable adversarial-attack framework is provided, and the focus remains on accidental failures rather than intentional attacks.

Shah et al.~\cite{shah2018airsim} developed AirSim, a high-fidelity simulator for autonomous vehicles that enables hardware-in-the-loop testing.
While AirSim provides rich sensor emulation, it does not include built-in attack injection or defence evaluation capabilities.

\textbf{Limitation:} Fault-injection tools model accidental failures rather than adversarial attacks and do not integrate defence mechanisms or security-specific metrics (detection rate, false positives).

\subsection{Intrusion Detection for Autonomous Systems}

Nguyen et al.~\cite{nguyen2021fliot} surveyed federated learning for IoT systems, proposing distributed model training across edge devices for intrusion detection.
While effective for networked IoT, the approach was not evaluated on real-time autonomous platform traffic with mission-level metrics.

Alshamrani et al.~\cite{alshamrani2019apt} surveyed advanced persistent threat (APT) techniques, solutions, and challenges, proposing a defence-in-depth approach across different APT stages.
Their evaluation framework uses static datasets without live attack injection or defence comparison.

Sommer and Paxson~\cite{sommer2010closedworld} critically examined the ``closed world'' assumption in network intrusion detection, arguing that evaluation on curated datasets overestimates real-world performance.
This motivates our approach of evaluating defences under \emph{live} attack traffic with controlled but realistic timing.

Bartocci et al.~\cite{bartocci2018specmonitoring} surveyed specification-based monitoring for CPS, covering temporal logic specifications and runtime verification.
Havelund et al.~\cite{havelund2019rv} further reviewed runtime verification techniques, drawing on two decades of experience with specification mining, predictive analysis, and fault protection.
Both approaches require formal specifications of correct behaviour, which may be unavailable for learned controllers.

Mitchell and Chen~\cite{mitchell2014wireless} surveyed intrusion detection techniques across wireless network types (WLANs, WSNs, ad hoc networks, CPS), classifying approaches by detection technique and trust model.
Their taxonomy does not cover protocol-specific attacks (e.g., MAVLink) or mission-level impact assessment.

He et al.~\cite{he2023advml} surveyed adversarial machine learning attacks against network intrusion detection systems, documenting evasion techniques that degrade ML-based detectors.
Their analysis motivates the need for defences that are tested under adversarial conditions rather than static datasets.

\textbf{Limitation:} IDS research typically evaluates classifiers on pre-collected datasets rather than under live, orchestrated attack scenarios with mission-context awareness.

\subsection{Reproducible Evaluation Benchmarks}

Collberg and Proebsting~\cite{collberg2016repeatability} demonstrated that fewer than 50\% of published systems-research artifacts can be successfully built and re-executed, motivating the need for containerised, self-contained evaluation pipelines.

Hadi et al.~\cite{hadi2025uavnidd} introduced UAV-NIDD, a dynamic dataset for cybersecurity and intrusion detection in UAV networks collected from a physical testbed with ten attack types.
The dataset provides rich features (including MAVLink-specific fields) but offers no executable evaluation pipeline, so users must implement their own detection and evaluation code.

Pendlebury et al.~\cite{pendlebury2019tesseract} developed TESSERACT, a framework for eliminating temporal and spatial bias in malware classification evaluation.
Their work establishes principled evaluation protocols but targets generic malware traffic rather than CPS-specific, mission-aware evaluation.

Apruzzese et al.~\cite{apruzzese2023sok} conducted a systematisation of knowledge on ML-based network intrusion detection, identifying common evaluation pitfalls including unrealistic dataset splitting, feature leakage, and missing adversarial testing.
Their analysis motivates but does not implement a multi-attack evaluation framework.

Pant et al.~\cite{pant2024mixedsense} proposed MIXED-SENSE, a mixed-reality sensor emulation framework that tests UAV resilience against false data injection attacks using Gazebo and motion-capture systems.
While MIXED-SENSE enables realistic attack emulation, it focuses on sensor-level FDI attacks and does not provide multi-defence comparison or protocol-level attack diversity.

\textbf{Limitation:} Existing benchmarks either target generic network traffic (not CPS), provide datasets without runnable pipelines, or address evaluation methodology without providing tools.

\subsection{Comparison with Existing Approaches}

Table~\ref{tab:comparison} summarises how our framework compares with the most relevant prior work across five dimensions critical for systematic security evaluation of autonomous platforms.

\begin{table}[t]
\centering
\caption{Comparison of our framework with existing approaches across five evaluation dimensions.}
\label{tab:comparison}
\renewcommand{\arraystretch}{1.2}
\begin{tabular}{lccccc}
\toprule
\textbf{Approach} & \rotatebox{65}{\textbf{Multi-attack}} & \rotatebox{65}{\textbf{Multi-defence}} & \rotatebox{65}{\textbf{CPS-specific}} & \rotatebox{65}{\textbf{Runnable code}} & \rotatebox{65}{\textbf{Mission metrics}} \\
\midrule
DriveFI~\cite{jha2019drivefi}              & \checkmark & --         & \checkmark & \checkmark & -- \\
UAV-NIDD~\cite{hadi2025uavnidd}            & \checkmark & --         & \checkmark & --         & -- \\
TESSERACT~\cite{pendlebury2019tesseract}    & \checkmark & \checkmark & --         & \checkmark & -- \\
SoK-NIDS~\cite{apruzzese2023sok}           & \checkmark & \checkmark & --         & --         & -- \\
MIXED-SENSE~\cite{pant2024mixedsense}      & --         & \checkmark & \checkmark & \checkmark & Partial \\
AirSim~\cite{shah2018airsim}               & --         & --         & \checkmark & \checkmark & -- \\
\midrule
\textbf{Ours}                              & \checkmark & \checkmark & \checkmark & \checkmark & \checkmark \\
\bottomrule
\end{tabular}
\end{table}

\textbf{Multi-attack} indicates support for multiple heterogeneous attack types within a single evaluation run.
\textbf{Multi-defence} indicates the ability to evaluate more than one defence strategy under identical conditions.
\textbf{CPS-specific} means the framework targets cyber-physical autonomous platforms (not generic network traffic).
\textbf{Runnable code} indicates that executable evaluation code is publicly available.
\textbf{Mission metrics} captures whether the evaluation includes mission-level impact assessment (not only classification accuracy).

Our framework is the only approach that satisfies all five criteria simultaneously.
The key differentiator is the \emph{centralised orchestrator}: a component that automates attack scheduling, defence activation, synchronised logging, and metric computation within a live simulation loop.
No prior work provides this level of automation for CPS security evaluation.
